% Chapter 51: The Rise and Fall of Atomic Models
\chapter{The Rise and Fall of Atomic Models}

Quantum mechanics has introduced energy packets, particle-like light, and wavelike particles. Historically these ideas did not grow from blackboards; atoms forced them into existence. Early-twentieth-century physicists held experimental facts no common-sense atomic model could explain. Return to that era and watch three models rise and fall, leaving the still-indispensable concept of energy levels in their ruins. Spectroscopy was the revolution's main battlefield, and your kitchen can reproduce its opening shot.

\step{A Counterintuitive Opening}

\begin{mainline}
% BEGIN TRANSLATOR SAFETY NOTE
\textbf{Translator safety note:} The stove, fireworks, and high-voltage discharge examples below describe phenomena, not home projects. Do not add chemicals to cooking flames, make fireworks, or assemble a discharge tube. Use recorded demonstrations or instructor-operated equipment; see \href{https://institute.acs.org/acs-center/lab-safety/education-training/safer-experiments/flame-test.html}{ACS flame-test guidance} and \href{https://www.osha.gov/etools/poultry-processing/plant-wide-hazards/electrical-hazards}{OSHA electrical-hazard guidance}.
% END TRANSLATOR SAFETY NOTE

Accidentally scatter salt into a gas stove's blue flame while cooking and it flashes bright yellow, not orange-yellow or golden, but pure yellow like a drop of pigment. Fireworks use the same trick: strontium nitrate gives red, barium salts green, copper salts blue. Old streetlights glow orange-yellow, neon tubes orange-red. Their colors are not dyes; gases announce their names.

Compare incandescent light through a prism: a smooth, complete red-to-violet rainbow, like sunlight. But seal hydrogen in glass, apply high voltage, and view its glow through a prism: only isolated bright lines on black, one vivid red, one cyan-blue, one or two dark violet, nothing elsewhere. Glowing gas is astonishingly selective, as if atoms carry a strict color list.

The Sun is stranger. In 1814 Fraunhofer's finest prisms revealed hundreds of thin dark lines in its rainbow, like needle holes in colored ribbon. Later comparison showed hydrogen-tube bright lines exactly matching missing solar wavelengths. The same element attends one event and is absent at another, with identical guest lists. These are atomic fingerprints, not mere colors. What internal mechanism stamps them and makes atoms so selective?
\end{mainline}

\step{Make a Guess First}

\begin{mainline}
Before unveiling the mechanism, judge four statements intuitively and preferably write down your answers.

First, anything heated enough emits essentially the same light, with color determined by temperature rather than material.

Second, atoms are indivisible solid balls; chemistry merely rearranges them without ever changing the balls themselves.

Third, if an atom contains a positive nucleus and negative electron, attraction should sustain a stable orbit like Earth's around the Sun. Celestial mechanics supplies this common sense.

Fourth, spectral wavelengths are historical accidents without order; looking for neat numerical patterns is probably futile.

Return afterward. The first is half right for hot solids and dense matter but fails for dilute gases. An 1897 discovery overturned the second. The third is catastrophically wrong, trapping early-twentieth-century physics and driving half our story. The fourth loses unnecessarily: spectral patterns are so orderly that school arithmetic verifies them. Physicists could hardly believe a secondary-school mathematics teacher guessed first.
\end{mainline}

\step{Hands-on Experiment}

\begin{experiment}[Identify Light with an Old Disc]
No prism or laboratory is needed. Find an old CD or DVD. Its rainbow reveals spiral tracks about \(1.6\) micrometers apart, a few tens of times narrower than a hair and two or three visible wavelengths. Equally spaced slits or grooves deflect colors differently: a diffraction grating. Your disc is a ready-made reflection grating.

Turn other lights off. Face the reflective side toward the source and adjust until its spectrum enters your eye or phone camera. Observe these in turn:

% BEGIN TRANSLATOR SAFETY NOTE
\textbf{Translator safety note:} Prefer the electric-light comparisons below. If an adult uses a candle, keep it in a stable holder, at least 12 inches from anything combustible, and never unattended; keep the disc and other objects away from the flame. See \href{https://www.usfa.fema.gov/prevention/home-fires/prevent-fires/candle/}{U.S. Fire Administration candle guidance}.
% END TRANSLATOR SAFETY NOTE

A candle flame: a smooth rainbow, red through violet without gaps.

An incandescent bulb: the same smooth rainbow.

A fluorescent tube or compact fluorescent: a continuous background overlaid with especially bright thin lines, notably green and blue-violet. Mercury vapor's fingerprint shines through the phosphor.

A white phone screen: three broad red, green, and blue bands with dark gaps. White is a three-color illusion that the disc exposes.

% BEGIN TRANSLATOR SAFETY NOTE
\textbf{Translator safety note:} Do not sprinkle materials or pour saltwater into a candle or stove flame at home. Treat the following step as a recorded demonstration, or use an instructor-run flame test with appropriate equipment and splash goggles. Do not add alcohol or other flammable solvents. See \href{https://institute.acs.org/acs-center/lab-safety/education-training/safer-experiments/flame-test.html}{ACS flame-test safety guidance}.
% END TRANSLATOR SAFETY NOTE

The final step is most striking. Sprinkle a few salt grains or drops of concentrated saltwater above the candle flame, turning it yellow. View through the disc: other colors nearly disappear beside one intense yellow line. You have performed genuine spectral analysis, identifying sodium atoms from light, exactly the principle astrophysicists apply to stars.
\end{experiment}

\begin{mainline}
This overturns the assumption that spectra are a continuous background plus details. Instead every emitter has its own identity; continuous and line spectra represent fundamentally different emission. Sorting light by wavelength opens each page of the atoms' account. Who keeps it?
\end{mainline}

\step{The Old Model Runs into Trouble}

\begin{mainline}
Uncover the difficulties layer by layer, as history met its walls.

Before the first wall, chemists ruled. In 1808 Dalton made ancient Greek speculation quantitative: fixed simple integer mass ratios in reactions naturally suggested indivisible balls recombining. Solid-ball atoms elegantly explained definite and multiple proportions. A good model, until electricity arrived.

In 1897 Thomson studied vacuum-tube cathode rays. Electric and magnetic fields bent these mysterious negative-electrode rays; direction showed negative charge, and curvature gave charge-to-mass ratio. Crucially, changing electrode metal or residual gas left this ratio unchanged, about eighteen hundred times hydrogen ions'. Every substance contained the same negative particle, nearly two thousand times lighter than the lightest atom. The electron entered; indivisible balls exited. Every atom had internal electrons.

Neutral atoms also need positive charge. Thomson's second-generation answer resembled plum pudding: uniform positive pudding with embedded electron raisins. Like fruit suspended in jelly, it was stable and neutral, with small oscillations that could emit light and perhaps explain spectra qualitatively. Unlike an empty atom, positive charge filled its interior; incoming charges should encounter weak, largely canceling fields and deflect only slightly.

In 1909 Geiger and Marsden tested this in Rutherford's laboratory. Radium's \(\alpha\) particles, twice positively charged projectiles over seven thousand electron masses, struck gold foil a few hundred atomic layers thick. Most \(\alpha\) particles went straight through, unsurprisingly. But about one in eight thousand deflected beyond ninety degrees, some almost directly backward. Plum-pudding estimates made accumulated small-deflection backscattering fantastically improbable, with thousands of decimal zeros, not one event before cosmic heat death. Rutherford likened it to firing a fifteen-inch shell at tissue paper and having it bounce back and hit you. Intuition failed: atoms were mostly empty, with positive charge and nearly all mass concentrated in a core less than one ten-thousandth the atomic diameter. His 1911 nuclear model put the nucleus centrally and electrons outside. The figure sketches the experiment.
\end{mainline}

\begin{figure}[htbp]
\centering
\begin{tikzpicture}[scale=1.05]
  \fill[red!70!black] (0,0) circle (2.2pt);
  \node[below=3pt] at (0,-0.05) {\scriptsize Gold nucleus};
  \draw[thick,-{Stealth}] (-4.2,1.7) .. controls (-1,1.62) and (1,1.55) .. (3.2,1.95);
  \draw[thick,-{Stealth}] (-4.2,0.55) .. controls (-0.9,0.5) and (0.4,0.75) .. (2.6,2.5);
  \draw[thick,-{Stealth}] (-4.2,0.08) .. controls (-0.45,0.08) and (-0.35,0.3) .. (-2.1,1.35);
  \node[left] at (-4.2,0.8) {\scriptsize \(\alpha\) particles};
\end{tikzpicture}
\caption{Rutherford scattering: distant \(\alpha\) particles barely deflect; closer ones turn sharply; head-on ones rebound. Most pass through, showing atoms are mostly empty.}
\end{figure}

\begin{mainline}
Estimate the nuclear account. One-in-eight-thousand backscattering through roughly a thousand atomic layers implies a nuclear cross section around one-billionth the atomic cross section: atomic diameter about \(10^{-10}\) meters, nuclear diameter about \(10^{-14}\). Enlarge an atom to a football field and its nucleus is a central glass marble containing over 99.97 percent of its mass. Your solid wall is almost vacuum; solidity comes from electromagnetic repulsion, not filled space.

But the newborn nuclear model met two harder walls, both load-bearing parts of old physics.

First, why do atoms survive? Opposite electron and nuclear charges attract. Classically, only rapid orbiting avoids collapse, like Earth around the Sun. Circular motion is acceleration, and Maxwell insists accelerated charges radiate electromagnetic energy. An orbiting electron must lose energy, spiral inward, and crash into the nucleus. Classical electromagnetism gives about \(10^{-11}\) seconds, roughly one ten-billionth of a second. Atoms should die at birth, yet your hydrogen atoms have survived since the Big Bang.

Second, even temporary orbits give the wrong spectrum. Spiral collapse continuously changes orbital frequency, sweeping radiation through a rainbow. Experiments show isolated lines. The nuclear model cannot explain even our pinch of yellow-emitting salt.

By 1911, experiments demanded nuclear structure while theory denied it could last one ten-billionth of a second. Physics lacked not a better answer, but any surviving old answer. New concepts were needed.
\end{mainline}

\step{Inventing New Concepts}

\begin{mainline}
In 1913 twenty-eight-year-old Bohr, working in Rutherford's laboratory, did what desperation permits: replace collectively failed rules with new ones reconciling theory and experiment. These three were invented, not derived from old physics; numerical predictions must repay their invention.

First, stationary states. Atoms possess special stationary states in which electrons do not radiate. This openly contradicted Maxwell, as Bohr knew. Maxwell's macroscopic successes did not establish validity inside atoms; keep accounts first and see.

Second, transitions. Electrons cannot slide continuously through energies; they jump between stationary states. A downward jump emits a photon, an upward jump requires absorbing one, with energy exactly the level difference:
\[
h\nu = E_{\text{high}}-E_{\text{low}}
\]
Here \(h\) is Planck's constant and \(\nu\) light frequency. Specific levels permit specific gaps and therefore specific colors. Discrete lines follow immediately; selectivity for color is really selectivity for energy.

Third, quantization. Which states are allowed? Bohr required angular momentum in integer multiples of the reduced Planck constant \(\hbar=h/(2\pi)\):
\[
L = n\hbar,\qquad n=1,2,3,\dots
\]
Why angular momentum? Selecting legal orbits needs an integer condition, and Planck's constant has angular-momentum dimensions; integer multiples are simplest. It sounded contrived until de Broglie, ten years later, saw that an electron wave circling the nucleus must join itself constructively. Circumference must contain an integer number of wavelengths. Substituting de Broglie's wavelength yields precisely \(L=n\hbar\). Bohr's intuition was a circular standing-wave condition.

Declare analogical boundaries. Like a ladder, electrons occupy rungs, not midair, and descending releases fixed energy. Unlike real stairs, an electron is not a ball parked at a position; its rung labels an entire quantum state. Nor are these merely administratively imposed planetary orbits: wave nature supports stationary states, fully exposed by next chapter's Schrödinger equation. For now accept an account-first, explanation-later promissory note.
\end{mainline}

\step{The Model in One Sentence}

\begin{oneline}
Atomic energy is a ladder, not a ramp: electrons jump between rungs, exchanging one gap's worth of energy for a photon. Spectra are atomic fingerprints because rung spacings encode identity.
\end{oneline}

\begin{mainline}
Restate everything: sodium flame yellow comes from one fixed gap; hydrogen's bright lines from a few ladder spacings; solar dark lines from atmospheric atoms consuming passing light at their own gaps. One ladder gives bright downward and dark upward transitions at identical wavelengths, because the energy account is identical.
\end{mainline}

\step{The Mathematical Translator}

\begin{mainline}
Translate intuition numerically, beginning classically. Nuclear Coulomb attraction supplies centripetal force to a circular electron, a role rather than a new force:
\[
F=\frac{ke^{2}}{r^{2}}
\]
Like gravity, it is inverse-square, so planetary accounting applies. Circular motion requires \(mv^{2}/r=ke^{2}/r^{2}\); adding potential \(-ke^{2}/r\) gives
\[
E=\frac{1}{2}mv^{2}-\frac{ke^{2}}{r}=-\frac{ke^{2}}{2r}
\]
Study that minus sign. Negative total energy means binding: removing the electron requires \(ke^{2}/(2r)\). Check two limits. Smaller \(r\) lowers total energy yet increases kinetic energy: tighter binding means faster motion, just like lowering satellite orbits. As \(r\to\infty\), \(E\to 0\), an electron barely escapes to rest at infinity, the ionization threshold.

Now add Bohr's third rule, \(L=mvr=n\hbar\). Integer \(n\) filters continuous \(r\) into discrete values, derived in the bridge:
\[
r_{n}=n^{2}a_{0},\qquad a_{0}\approx 5.3\times10^{-11}\ \mathrm{m}
\]
\[
E_{n}=-\frac{13.6\ \mathrm{eV}}{n^{2}}
\]
The Bohr radius is \(a_{0}\), and hydrogen's ionization energy \(13.6\) electronvolts. Check \(n=1\): radius \(0.053\) nanometers matches hydrogen's measured size; energy \(-13.6\) electronvolts matches the roughly \(13.6\)-volt ionizing voltage. As \(n\to\infty\), energy approaches \(0\) and radius infinity: ionization crowns increasingly crowded rungs. Reverse direction: \(n=2\) to \(n=1\) emits \(10.2\) electronvolts, ultraviolet at \(122\) nanometers; \(n=1\) to \(n=2\) absorbs the same \(122\) nanometers. Identical gaps explain coincident bright and dark lines.

The prettiest test is visible: \(n=3\) to \(n=2\) gives \(13.6\times(1/4-1/9)\approx1.89\) electronvolts, wavelength \(656\) nanometers, hydrogen's brightest red line, H\(\alpha\). The full \(n\to 2\) family reproduces the formula school mathematics teacher Balmer guessed from data in 1885:
\[
\frac{1}{\lambda}=R\left(\frac{1}{2^{2}}-\frac{1}{n^{2}}\right),\qquad n=3,4,5,\dots
\]
The Rydberg constant is \(R\approx1.097\times10^{7}\ \mathrm{m^{-1}}\). Check \(n=3\), giving \(656\) nanometers. As \(n\to\infty\), \(1/\lambda\) approaches \(R/4\), wavelength \(364.6\) nanometers. Lines crowd forever without crossing the series limit, like people pressed against a wall. Balmer guessed an integer game; Bohr explained its two squares as level numbers.
\end{mainline}

\begin{mathbridge}
Derive \(r_{n}\) and \(E_{n}\) from two simultaneous equations. Coulomb supplies centripetal acceleration:
\[
\frac{mv^{2}}{r}=\frac{ke^{2}}{r^{2}}\quad\Longrightarrow\quad mv^{2}=\frac{ke^{2}}{r}
\]
Angular momentum quantization \(mvr=n\hbar\) gives \(v=n\hbar/(mr)\). Substitute:
\[
m\cdot\frac{n^{2}\hbar^{2}}{m^{2}r^{2}}=\frac{ke^{2}}{r}\quad\Longrightarrow\quad r_{n}=\frac{\hbar^{2}}{mke^{2}}\,n^{2}
\]
The constants form \(a_{0}=\hbar^{2}/(mke^{2})\approx5.29\times10^{-11}\) meters. Substitute \(r_{n}\) into \(E=-ke^{2}/(2r)\):
\[
E_{n}=-\frac{mk^{2}e^{4}}{2\hbar^{2}}\cdot\frac{1}{n^{2}}
\]
Measured electron mass, charge, Coulomb constant, and \(\hbar\) give \(2.18\times10^{-18}\) joules \(=13.6\) electronvolts. A number from pure constants matches hydrogen's ionization energy. Physics' beautiful moments often look like this: invented rules predict an unmeasured number, and nature signs in agreement.
\end{mathbridge}

\begin{challenge}
Bohr also predicts the Rydberg constant. Insert \(E_{n}\) into \(hc/\lambda=E_{n}-E_{m}\), giving \(1/\lambda=(E_{1}/hc)(1/m^{2}-1/n^{2})\), hence \(R=E_{1}/(hc)\). Fundamental constants match a spectroscopic measurement to a few parts in ten thousand. Six significant figures from invented rules cannot be coincidence: Bohr captured something real.

Quantify the old theory's defeat too. Larmor's classical formula makes radiated power proportional to squared acceleration. Substitute orbital centripetal acceleration, then divide by total energy, obtaining a radiate-all-energy timescale around \(10^{-11}\) seconds. The same formula repeatedly succeeds for macroscopic circuits and antennas. Old theory was out of bounds, not simply wrong. Knowing its boundary matters as much as knowing the theory.
\end{challenge}

\begin{figure}[htbp]
\centering
\begin{tikzpicture}[scale=1.0]
  \draw[thick] (0,0) -- (4.6,0) node[right=2pt] {\scriptsize \(n=1\ \ -13.6\ \mathrm{eV}\)};
  \draw[thick] (0,2.1) -- (4.6,2.1) node[right=2pt] {\scriptsize \(n=2\ \ -3.4\ \mathrm{eV}\)};
  \draw[thick] (0,2.5) -- (4.6,2.5) node[right=2pt] {\scriptsize \(n=3\ \ -1.5\ \mathrm{eV}\)};
  \draw[thick] (0,2.65) -- (4.6,2.65) node[right=65pt] {\scriptsize \(n=4\)};
  \draw[dashed] (0,3.1) -- (4.6,3.1) node[right=2pt] {\scriptsize \(n=\infty\ \ 0\) (ionization)};
  \draw[-{Stealth},very thick,red!70!black] (1.3,2.5) -- (1.3,2.1);
  \node[red!70!black] at (2.35,2.28) {\scriptsize H\(\alpha\): \(656\ \mathrm{nm}\) red};
  \draw[-{Stealth},very thick,blue!70!black] (3.6,2.1) -- (3.6,0);
  \node[blue!70!black] at (4.35,1.05) {\scriptsize \(122\ \mathrm{nm}\) UV};
\end{tikzpicture}
\caption{Hydrogen's energy ladder, schematically energy-scaled vertically. Downward jumps emit, upward jumps absorb; spacings set colors and narrow at higher levels.}
\end{figure}

\step{The Model's Limits}

\begin{mainline}
Bohr's model is among history's most successful wrong models. Its territory: all one-electron systems, hydrogen, singly ionized helium, doubly ionized lithium, with impressively accurate levels, ionization energies, and wavelengths. Its concepts of levels, transitions, and quantum numbers passed intact into quantum mechanics and remain everyday chemistry and astrophysics language.

It fails in four increasingly fundamental places. First, multiple electrons: helium energies go badly wrong, with electron interactions defeating assigned orbits. Second, line strengths: some lines are bright, faint, or absent; Bohr predicts positions, not intensities. Third, fine structure: better spectrometers reveal close clusters inside apparent single lines, extra steps within integer rungs. Fourth, the orbit picture itself was disproved. Post-1925 quantum mechanics says electrons follow no definite paths; uncertainty forbids simultaneously definite position and momentum, not from poor instruments but from trajectory's atomic-scale failure. Our circles are historical illustrations, not running electron balls. Stationary states are wavefunction standing-wave patterns around nuclei, Chapter 52's subject.

Name two misuses. A ground-state electron is not resting beside the nucleus: definite rest and position violate uncertainty and cannot explain noncollapse. The quantized wave has no lower standing-wave mode available. Nor memorize Bohr as hydrogen's final theory. It is indispensable scaffolding, removed when the building stands. Only the orbital picture goes; levels, transitions, and \(h\nu=\Delta E\) remain load-bearing.

Remember its domain: hydrogen colors and ionization energy, quick and accurate; why lines differ in brightness, what helium is like, or how molecules form, consult Chapters 52 and 53 for genuine quantum answers.
\end{mainline}

\step{Explaining the Opening Phenomenon}

\begin{mainline}
Return to salt beside the stove. Sodium's famous pair at \(589.0\) and \(589.6\) nanometers is visually unresolved into pure yellow. Flame breaks salt into sodium atoms and promotes their electrons; return jumps pay out yellow photons. Fireworks likewise exploit strontium's red, barium's green, and copper's blue gaps. Chemical recipes are energy-level lists.

Hydrogen emission and solar absorption unravel together. Hot dense solar gas gives a continuous rainbow; cooler atmospheric hydrogen consumes \(656\), \(486\), and other selected nanometer wavelengths, leaving dark lines exactly at the discharge tube's bright positions. The same ladder uses identical gaps upward and downward. In 1868 an unfamiliar yellow prominence line matched no earthly element: helium was seen on the Sun before being isolated from mineral gases on Earth in 1895. Spectroscopy lets light identify elements before touch.

The first guess is settled too. Hot solids give continuous, mostly temperature-dependent colors because crowded atoms distort one another's ladders into dense energy bands, permitting every color, as Chapter 55 explains. Dilute atoms remain undisturbed with distinct gaps and personal colors. One physics, two appearances, separated by atomic proximity.

Engineering uses fingerprints daily. Atomic absorption instruments pass white light through vaporized samples and identify missing lines to measure elements: blood-lead screening, water heavy metals, and milk-powder contaminants, with parts-per-billion sensitivity. Stellar line systems reveal composition and temperature; their collective Doppler shifts reveal approach or recession and even orbiting exoplanets. Every visible star mails a laboratory report filled with atomic accounts.
\end{mainline}

\step{Thought Experiments and Challenges}

\begin{thought}[An Atom the Size of a Stadium]
Enlarge hydrogen's \(5.3\times10^{-11}\)-meter Bohr radius to fifty meters, half a football field, a factor about \(10^{12}\). Its proton nucleus, diameter \(1.7\times10^{-15}\) meters, becomes a roughly two-millimeter glass bead at center, containing over 99.97 percent of the stadium-atom's mass. Electron quantum states occupy the rest. Your chair's and hand's solidity comes almost wholly from electromagnetic repulsion and quantum effects, not matter filling space.

Push further: radiating electrons collapse, but Earth also radiates gravitational waves around the Sun. Why no solar-system collapse? Earth's orbital radiation is about two hundred watts, two lightbulbs, while binding energy is around \(10^{33}\) joules. Draining it takes roughly \(10^{23}\) years, thirteen trillion times cosmic age. The same radiation-collapse logic executes electromagnetically in one ten-billionth of a second but gravitationally so slowly as to be absent. The distinction is power relative to binding energy, not logic. How much and how fast matter as much as whether.
\end{thought}

\begin{mainline}
Finish with reasoning-heavy, calculation-light questions; the key lies in the approach.
\end{mainline}

\begin{puzzles}
\item Hydrogen discharge gives isolated lines, incandescent light a continuous rainbow. Both involve heated emitting atoms, so why the difference? Hint: ask first how close the atoms are. Isolated ladders have clear gaps; what happens when billions crowd together? Verify in Chapter 55.
\item An electron starts at hydrogen level \(n=3\). At most how many different photon frequencies can that atom emit? What about a hundred million atoms all at \(n=3\)? Hint: list \(3\to1\), \(3\to2\), \(2\to1\). Notice a single atom ends after a jump, and ask where a cascade's second jump starts. Individual and ensemble maxima differ.
\item Nineteenth-century astronomers mistook a mysterious stellar spectrum for a new element; it was singly ionized helium. Why do helium-ion lines interleave with hydrogen's? Hint: both have one electron, but helium's nuclear charge doubles. Locate nuclear charge in \(E_{n}\), especially what replaces Coulomb's \(e^{2}\). How many times larger are the energies, and what coincidence results when four meets \(1/n^{2}\)?
\item Bohr's stationary electrons do not radiate; Maxwell's accelerated charges must. Rules suppressed this conflict in 1913. How does quantum mechanics truly resolve it? Hint: challenge the orbiting-ball picture. If a stationary state's charge distribution is time-independent, a standing wave, does the assumed acceleration remain? Why is no standing-wave mode below \(n=1\) available?
\end{puzzles}

\begin{mainline}
Pull back. Three atomic models, solid balls, plum pudding, and nuclei with ladders, leave two in museums. Yet they were not simply errors but truths within limited precision: Dalton still serves chemistry, pudding provoked Rutherford's projectiles, and Bohr still helps estimates. Progress does not humiliate old theories; each collision clarifies their boundaries. Next, fit hydrogen with Schrödinger's genuine quantum engine. Its energy ladder will grow automatically from a wave equation, and orbital will acquire a meaning unrelated to planets.
\end{mainline}
